\chapter{Swaps de change, change à terme et base de devises}\label{ch:m2:fx-swaps-forwards-and-the-cross-currency-basis}

À la fin d'août 2026, un \omterm{def:m2:money-markets:bill}{bon du Trésor} américain à dix ans rapportait 4.75~\% et une
obligation d'État japonaise à dix ans, 2.94~\%. Un assureur japonais qui achète ce bon
gagne 1.81 point de pourcentage de plus, mais il doit à ses assurés des yens, et une baisse
du dollar effacerait ce gain plusieurs fois, alors il se couvre~: il vend les
dollars à terme, par tranches de trois mois. Avec le taux overnight du dollar à
3.68~\% et celui du yen inférieur à 1~\%, chaque roll de la couverture coûte environ 2.70
points de pourcentage par an, soit plus que le rendement supplémentaire, avant que le marché
ne facture rien de plus pour prêter des dollars. Le Trésor couvert rapporte environ 2.05~\%,
moins que l'obligation japonaise. Ce chapitre porte sur les instruments qui fixent
ce coût~: le \omterm{def:m2:fx-swaps-forwards-and-the-cross-currency-basis:forward}{change à terme sec}, le \omterm{def:m2:fx-swaps-forwards-and-the-cross-currency-basis:swap}{swap de change}, qui est de loin l'instrument le plus
échangé sur le marché des changes, la parité qui les relie aux taux d'intérêt,
et la \omterm{def:m2:fx-swaps-forwards-and-the-cross-currency-basis:basis}{base de devises}, par laquelle cette parité ne tient plus depuis 2007.

\section{Change à terme sec et points de terme}

\begin{definition}[Change à terme sec, points de terme]\label{def:m2:fx-swaps-forwards-and-the-cross-currency-basis:forward}
Un \emph{change à terme sec}\index{change à terme sec} (outright forward) est un accord pour échanger
deux devises à un taux $F$ fixé aujourd'hui, pour une valeur à une date postérieure au spot. Ses
\emph{points de terme}\index{points de terme} sont $F - S$ exprimés en pips de la
paire~: le marché cote les taux à terme sous forme de points ajoutés au taux spot.
\end{definition}

Un change à terme se fixe par un raisonnement qui ne fait appel qu'à des dépôts. Emprunter une
unité de la devise de base pour la période à son taux $r_b$, la vendre au spot contre $S$
unités de la devise de cotation, et prêter celles-ci au taux $r_q$ de cette devise~:
à l'échéance on doit $1 + r_b\tau_b$ de la devise de base et on détient $S(1 + r_q\tau_q)$ de la
devise de cotation, où les $\tau$ sont les fractions d'accumulation selon la base de décompte de chaque devise. Le
terme qui rend nul l'aller-retour est le change à terme de parité couverte.

\begin{definition}[Parité des taux d'intérêt couverte]\label{def:m2:fx-swaps-forwards-and-the-cross-currency-basis:cip}
La \emph{parité des taux d'intérêt couverte}\index{parité des taux d'intérêt couverte} (covered interest parity, CIP) est la
relation
\begin{equation}\label{eq:m2:fx-swaps-forwards-and-the-cross-currency-basis:cip}
F \;=\; S\,\frac{1 + r_q\,\tau_q}{1 + r_b\,\tau_b}
\end{equation}
entre le spot, le terme et les taux du marché monétaire des deux devises sur la
même période.
\end{definition}

La devise au taux d'intérêt le plus élevé se traite à terme avec un déport~: son
terme achète moins de l'autre devise que son spot ne le fait, et le rendement plus faible de
l'autre devise est compensé par le gain de change. USDJPY est coté en yens par
dollar~; avec des taux en dollars supérieurs à ceux en yens, ses \omterm{def:m2:fx-swaps-forwards-and-the-cross-currency-basis:forward}{points de terme} sont
négatifs.

\begin{example}[USDJPY à trois mois]\label{ex:m2:fx-swaps-forwards-and-the-cross-currency-basis:fwd}
Avec USDJPY à 156.87, un taux en dollars de 3.68~\% (actual/360) et un taux en yens de
0.977~\% (actual/365), le terme à 91 jours donné par l'\cref{eq:m2:fx-swaps-forwards-and-the-cross-currency-basis:cip}
vaut 155.8028~: \omterm{def:m2:fx-swaps-forwards-and-the-cross-currency-basis:forward}{points de terme} de $-106.7$ pips de 0.01 yen. Une \omterm{def:m2:fx-swaps-forwards-and-the-cross-currency-basis:basis}{base de devises} de $-25$ points
de base sur la jambe en yens, telle que définie plus bas, la ramène à 155.7059, $-116.4$ pips.
Les points \omterm{def:m2:fx-swaps-forwards-and-the-cross-currency-basis:swap}{tom-next} à un jour valent $-1.18$.
\end{example}

\section{Swaps de change et roll à court terme}

\begin{definition}[Swap de change, tom-next]\label{def:m2:fx-swaps-forwards-and-the-cross-currency-basis:swap}
Un \emph{swap de change}\index{swap de change} (FX swap) est une paire d'échanges opposés portant sur le même
montant d'une devise~: une jambe proche, en général au spot, et une jambe lointaine à une date à terme,
les deux taux différant des \omterm{def:m2:fx-swaps-forwards-and-the-cross-currency-basis:forward}{points de terme}. Le \emph{tom-next}\index{tom-next}
est le swap d'un jour allant du prochain jour ouvré (demain) au jour d'après
(spot next), qui sert à reporter les positions d'une date de spot à la suivante.
\end{definition}

Un \omterm{def:m2:fx-swaps-forwards-and-the-cross-currency-basis:swap}{swap de change} est un prêt collatéralisé. La partie qui achète des dollars à la
jambe proche et les revend à la jambe lointaine a emprunté des dollars et prêté l'autre
devise pour la période, chaque prêt étant garanti par l'autre~; les \omterm{def:m2:fx-swaps-forwards-and-the-cross-currency-basis:forward}{points de terme} sont
l'écart de taux d'intérêt. C'est de loin le plus important des instruments de change~:
4~billions de dollars par jour en avril 2025, soit 42~\% de l'ensemble des échanges de change, surtout pour sept jours ou
moins (traité au \cref{ch:m2:the-fx-market}). Les banques l'utilisent pour financer leurs actifs en dollars,
les investisseurs pour couvrir leurs avoirs à l'étranger, et un trader qui détient une position au
spot la reporte chaque jour au moyen d'un swap \omterm{def:m2:fx-swaps-forwards-and-the-cross-currency-basis:swap}{tom-next}, en payant ou en encaissant les points du
jour~: c'est le carry d'une position de change.

\begin{omfigure}
\centering
\begin{tikzpicture}[font=\small,
  box/.style={draw,thick,rounded corners=2pt,align=center,minimum height=1cm,minimum width=2.8cm,inner sep=3pt},
  arr/.style={-{Stealth[length=2.5mm]},thick}]
\node[box,draw=omDef,fill=omDef!8] (inv) at (0,0) {Investisseur japonais};
\node[box,draw=omThm,fill=omThm!8] (bk) at (9.2,0) {prêteur en dollars\\(banque)};
\node[font=\footnotesize\bfseries,anchor=west] at (-1.4,0.95) {jambe proche, spot};
\node[font=\footnotesize\bfseries,anchor=west] at (-1.4,-1.05) {jambe lointaine, trois mois plus tard};
\draw[arr] ([yshift=5pt]inv.east) -- node[above,font=\footnotesize] {1\,568.70~millions de yens} ([yshift=5pt]bk.west);
\draw[arr] ([yshift=-5pt]bk.west) -- node[below,font=\footnotesize] {10~millions de dollars} ([yshift=-5pt]inv.east);
\draw[arr,dashed] ([yshift=-1.8cm]bk.west) -- node[above,font=\footnotesize] {1\,557.06~millions de yens} ([yshift=-1.8cm]inv.east);
\draw[arr,dashed] ([yshift=-2.3cm]inv.east) -- node[below,font=\footnotesize] {10~millions de dollars} ([yshift=-2.3cm]bk.west);
\end{tikzpicture}

\omcaption{Un \omterm{def:m2:fx-swaps-forwards-and-the-cross-currency-basis:swap}{swap de change} USDJPY à trois mois au terme donné dans
l'\cref{ex:m2:fx-swaps-forwards-and-the-cross-currency-basis:fwd} avec une base de $-25$
points de base. L'investisseur emprunte des dollars et prête des yens, chaque prêt étant garanti
par l'autre~; il récupère 11.64~millions de yens de moins que ce qu'il a payé, soit le prix des
dollars. À titre indicatif.}\label{fig:m2:fx-swaps-forwards-and-the-cross-currency-basis:legs}
\end{omfigure}

\section{Parité des taux d'intérêt couverte}

Avant 2007, le CIP tenait aux coûts de transaction près, car tout écart constituait un profit
sans risque pour une banque~: si les dollars étaient moins chers à emprunter par \omterm{def:m2:fx-swaps-forwards-and-the-cross-currency-basis:swap}{swaps de change}
que sur le marché monétaire, une banque les emprunterait par swaps et les prêterait au comptant,
et ses transactions refermeraient l'écart. Les économistes de la BIS l'ont appelé «~la
chose la plus proche d'une loi physique en finance internationale~». Les mêmes
économistes ont constaté en 2016 qu'elle avait été systématiquement violée depuis
la crise financière de 2007--2008.

\begin{proposition}[Points de terme par échéance]\label{prop:m2:fx-swaps-forwards-and-the-cross-currency-basis:points}
Sous le CIP avec des taux constants, les \omterm{def:m2:fx-swaps-forwards-and-the-cross-currency-basis:forward}{points de terme} d'une paire croissent presque en
proportion à l'échéance, $F - S \approx S\,(r_q - r_b)\,\tau$ pour les échéances courtes,
et une base $b$ sur la jambe en devise de cotation ajoute approximativement $S\,b\,\tau$.
\end{proposition}

\begin{proof}
On développe \cref{eq:m2:fx-swaps-forwards-and-the-cross-currency-basis:cip} au premier
ordre en $\tau$~: $F/S \approx 1 + r_q\tau - r_b\tau$. Avec $r_q + b$ à la place de
$r_q$, le terme $b\tau$ apparaît.
\end{proof}

\begin{omfigure}
\begin{tikzpicture}
\begin{axis}[width=11cm,height=4.8cm,
  xlabel={jours jusqu'à la jambe éloignée},ylabel={points de terme USDJPY},
  tick label style={font=\small},label style={font=\small},
  xmin=0,xmax=370,ymin=-480,ymax=10,grid=major,grid style={gray!25},
  legend columns=2,legend style={font=\footnotesize,at={(0.5,-0.32)},anchor=north,draw=none,fill=none,/tikz/every even column/.append style={column sep=8pt}},
  table/col sep=comma]
\addplot[omDef,very thick,mark=*,mark size=1.3pt] table[x=days,y=nobasis]{figdata/markets-2/16-fx-swaps-forwards-and-the-cross-currency-basis/points.csv};
\addplot[omThm,very thick,dashed,mark=*,mark size=1.3pt] table[x=days,y=basis]{figdata/markets-2/16-fx-swaps-forwards-and-the-cross-currency-basis/points.csv};
\legend{parité des taux d'intérêt couverte,avec une base de $-25$ pb sur la jambe en yen}
\end{axis}
\end{tikzpicture}

\omcaption{\omterm{def:m2:fx-swaps-forwards-and-the-cross-currency-basis:forward}{Points de terme} USDJPY, en pips de 0.01 yen, d'une semaine à un an,
avec des \omterm{def:m2:money-markets:yields}{taux monétaires} constants de 3.68~\% en dollars et de 0.977~\% en yen, pour un cours au comptant de
156.87. Les points croissent avec l'échéance~; une base négative sur le yen les rend
plus négatives, le dollar devenant plus cher à terme. À titre indicatif~; données~: le
tutoriel du chapitre.}\label{fig:m2:fx-swaps-forwards-and-the-cross-currency-basis:points}
\end{omfigure}

\section{La base de devises}

\begin{definition}[Swap de base de devises, base de devises]\label{def:m2:fx-swaps-forwards-and-the-cross-currency-basis:basis}
Un \emph{swap de base de devises}\index{swap de base de devises} échange
des notionnels dans deux devises au départ au cours au comptant, échange des intérêts
variables dans chaque devise pendant sa durée, l'une des jambes portant un spread, et
rééchange les mêmes notionnels à l'échéance. Le spread est la
\emph{base de devises}\index{base de devises} $b$, cotée sur la
jambe non dollar~; pour les échéances courtes, elle se mesure à partir des \omterm{def:m2:fx-swaps-forwards-and-the-cross-currency-basis:swap}{swaps de change} comme l'écart entre
le taux du yen (ou de l'euro) implicite dans les \omterm{def:m2:fx-swaps-forwards-and-the-cross-currency-basis:forward}{points de terme} et le taux au comptant~:
\[
b \;=\; \left(\frac{F}{S}\,(1 + r_b\tau_b) - 1\right)\frac{1}{\tau_q} - r_q.
\]
\end{definition}

Une base négative sur le yen signifie que prêter du yen et emprunter des dollars par
swaps rapporte moins que le taux du marché monétaire en yen~: les dollars coûtent plus cher
par swaps que sur le marché au comptant. Depuis 2007, la base contre le dollar est négative pour
la plupart des devises, notamment l'euro et le yen, et positive pour quelques-unes, comme
le dollar australien. L'explication de la BIS combine demande et limites. Demande~:
des investisseurs comme les assureurs japonais couvrent de larges avoirs en dollars en empruntant
des dollars par swaps, et les entreprises qui émettent des obligations en euros et échangent le produit
en dollars s'y ajoutent. Limites~: les banques qui arbitreraient l'écart paient désormais
pour le bilan que mobilise l'opération, en fonds propres et en ratio d'effet de levier, et ne le font
que pour un spread. Depuis 2014, la base bondit aux fins de trimestre, avec les \omterm{def:m2:repo-and-specials:rate}{taux de repo},
lorsque les banques réduisent leur bilan pour leurs états périodiques.

Pour l'investisseur de l'accroche, la base est un coût direct. Le roll d'une couverture
de trois mois sur une obligation en dollars coûte, sur un an, environ le taux en dollars moins le taux
en yen plus la base en yen prise avec le signe inversé~:
\begin{equation}\label{eq:m2:fx-swaps-forwards-and-the-cross-currency-basis:hedged}
y^{\mathrm{hedged}} \;\approx\; y_{\mathrm{USD}} - \bigl(r_{\mathrm{USD}} - r_{\mathrm{JPY}} - b\bigr).
\end{equation}

\begin{omfigure}
\begin{tikzpicture}
\begin{axis}[width=11cm,height=5cm,
  ylabel={pourcent},
  tick label style={font=\small},label style={font=\small},
  xmin=2018.2,xmax=2026.75,ymin=-2,ymax=5.5,grid=major,grid style={gray!25},
  xtick={2019,2020,2021,2022,2023,2024,2025,2026},xticklabel style={/pgf/number format/1000 sep={}},
  legend columns=3,legend style={font=\footnotesize,at={(0.5,-0.18)},anchor=north,draw=none,fill=none,/tikz/every even column/.append style={column sep=8pt}},
  table/col sep=comma]
\addplot[omDef,thick] table[x=t,y=ust]{figdata/markets-2/16-fx-swaps-forwards-and-the-cross-currency-basis/hedged.csv};
\addplot[omThm,very thick] table[x=t,y=hedged]{figdata/markets-2/16-fx-swaps-forwards-and-the-cross-currency-basis/hedged.csv};
\addplot[omProp,thick,dashed] table[x=t,y=jgb]{figdata/markets-2/16-fx-swaps-forwards-and-the-cross-currency-basis/hedged.csv};
\draw[gray] (axis cs:2018.2,0) -- (axis cs:2026.75,0);
\legend{Treasury à 10 ans,couvert en yen (base nulle),JGB à 10 ans}
\end{axis}
\end{tikzpicture}

\omcaption{Un Treasury à 10 ans pour un investisseur en yen, d'avril 2018 à août 2026~:
non couvert, couvert par le roll de \omterm{def:m2:fx-swaps-forwards-and-the-cross-currency-basis:swap}{swaps de change} à court terme sans base (le rendement du
Treasury moins le SOFR plus le taux du call japonais), et le JGB à 10 ans. Toute base négative
abaisse davantage la courbe couverte. Rendements du Treasury et SOFR en fin de mois, moyennes
mensuelles pour le Japon. Données~: séries FRED DGS10, SOFR, IRSTCI01JPM156N et
IRLTLT01JPM156N.}\label{fig:m2:fx-swaps-forwards-and-the-cross-currency-basis:hedged}
\end{omfigure}

\begin{dated}{2026-09}{Les données de l'accroche}\label{dat:m2:fx-swaps-forwards-and-the-cross-currency-basis:inputs}
Treasury à 10 ans à 4.75~\% et SOFR à 3.68~\% le 31 août 2026~; taux du call monétaire
japonais à 0.977~\% et JGB à 10 ans à 2.94~\%, en moyennes d'août 2026~; USDJPY à 156.87 le 18
septembre 2026 (Federal Reserve H.10). Les taux au jour le jour tiennent lieu de taux
à trois mois, et aucune série officielle de la base n'est utilisée~: ses niveaux dans ce chapitre
sont illustratifs.
\end{dated}

La \cref{fig:m2:fx-swaps-forwards-and-the-cross-currency-basis:hedged} montre
la conséquence. Lorsque les taux de la Federal Reserve étaient très supérieurs à ceux du Japon, de la
fin 2018 au début de 2020 et de nouveau à partir de novembre 2022, le Treasury couvert rapportait moins
que l'obligation japonaise, avant même toute base~: dans 57 des 101 mois représentés. L'investisseur
doit ensuite choisir entre un rendement couvert plus faible, une position non couverte qui parie sur le
yen, ou rester sur son marché national.

\section{Tensions sur le financement en dollars}

Les banques hors États-Unis détiennent de larges actifs en dollars, financés en partie par des \omterm{def:m2:fx-swaps-forwards-and-the-cross-currency-basis:swap}{swaps de change}. Lorsque
le financement en dollars se tarit, comme en 2008 et en mars 2020, la base s'élargit fortement,
et la banque centrale qui émet le dollar le prête à d'autres banques centrales,
qui le prêtent à leurs banques. La Federal Reserve a d'abord mis en place de telles lignes de swap
le 12 décembre 2007, avec la BCE jusqu'à concurrence de USD~20 milliards et la Banque
nationale suisse jusqu'à concurrence de USD~4 milliards~; le 13 octobre 2008, les lignes avec la Banque
d'Angleterre, la BCE et la BNS ont été relevées «~pour accueillir toute quantité de
financement en dollars américains demandée~». En mars 2020, les lignes permanentes avec cinq
banques centrales ont été ramenées au taux OIS du dollar plus 25 points de base le 15
mars, et des lignes temporaires avec neuf autres ont été ouvertes le 19 mars
(\cref{fig:m2:fx-swaps-forwards-and-the-cross-currency-basis:swaplines}).

\begin{omfigure}
\begin{tikzpicture}
\begin{axis}[width=11cm,height=4.6cm,
  ylabel={Md USD},
  tick label style={font=\small},label style={font=\small},
  xmin=2007,xmax=2026.75,ymin=0,ymax=650,grid=major,grid style={gray!25},
  xtick={2008,2010,2012,2014,2016,2018,2020,2022,2024,2026},xticklabel style={/pgf/number format/1000 sep={}},
  table/col sep=comma]
\addplot[omThm,very thick,const plot] table[x=t,y=bn]{figdata/markets-2/16-fx-swaps-forwards-and-the-cross-currency-basis/swaplines.csv};
\end{axis}
\end{tikzpicture}

\omcaption{Dollars prêtés par la Federal Reserve aux banques centrales étrangères au moyen
de lignes de swap, en fin de mois, de 2007 à août 2026. Les pics correspondent aux crises
de financement en dollars de la fin de 2008 (un maximum hebdomadaire de USD~583 milliards le 17 décembre 2008)
et du printemps 2020 (USD~449 milliards le 27 mai 2020). Données~: série FRED SWPT, à partir de la
publication H.4.1.}\label{fig:m2:fx-swaps-forwards-and-the-cross-currency-basis:swaplines}
\end{omfigure}

Les lignes de swap plafonnent la base en période de crise, puisqu'une banque ayant accès aux
enchères en dollars de sa banque centrale n'acceptera pas de payer beaucoup plus que leur taux.
Elles ne la suppriment pas en temps normal, lorsque les facteurs sont la demande de couverture
et le prix du bilan.

\section{Tutoriel~: calcul des points de terme et de la base implicite}\label{tut:m2:fx-swaps-forwards-and-the-cross-currency-basis:basis}

\begin{tutorial}
\textbf{Objectif.} Calculer les \omterm{def:m2:fx-swaps-forwards-and-the-cross-currency-basis:forward}{points de terme} à partir de deux taux du marché monétaire, en déduire
la base implicite des points de marché et calculer un rendement couvert.
\textbf{État final~:}
\cref{ex:m2:fx-swaps-forwards-and-the-cross-currency-basis:fwd},
\cref{fig:m2:fx-swaps-forwards-and-the-cross-currency-basis:points,fig:m2:fx-swaps-forwards-and-the-cross-currency-basis:hedged}
et les chiffres du problème du week-end.
\begin{tutsteps}
\item \textbf{Le terme et la base}, au moyen de
l'\cref{eq:m2:fx-swaps-forwards-and-the-cross-currency-basis:cip} et de son inverse.
\omcode{code/firm/fxfwd/firm_fxfwd.py}{10}{30}{Terme par parité couverte, points de terme, taux implicite et base.}\label{lst:m2:fx-swaps-forwards-and-the-cross-currency-basis:fwd}
\item \textbf{La couverture}, au moyen de
l'\cref{eq:m2:fx-swaps-forwards-and-the-cross-currency-basis:hedged}.
\omcode{code/firm/fxfwd/firm_fxfwd.py}{39}{47}{Coût de couverture et rendement couvert.}\label{lst:m2:fx-swaps-forwards-and-the-cross-currency-basis:hedge}
\item \textbf{Exécuter} \texttt{fxfwd\_demo.three\_month()},
\texttt{fxfwd\_demo.hedged\_now()} et \texttt{fig\_fxfwd.py}.
\end{tutsteps}
\textbf{Pour aller plus loin.} Remplacer les taux au jour le jour par des taux à terme de la
maturité de la couverture, et comparer le rendement couvert obtenu avec des swaps à trois mois
renouvelés par roll à celui d'un terme à un an.
\end{tutorial}

\section{Construction~: le calculateur de terme et de base}\label{bld:m2:fx-swaps-forwards-and-the-cross-currency-basis:fxfwd}

\begin{build}
\bfield{Rôle} L'entreprise miniature cote des termes et des swaps à ses clients, finance des positions
en dollars au moyen de swaps et couvre l'exposition de change des investisseurs~: elle a besoin de
termes issus de courbes, de la base à partir des points de marché et de rendements couverts. \bfield{Interface} \texttt{forward(spot, r\_quote, r\_base, days, basis\_quote,
basis\_base, dc\_quote, dc\_base)}~; \texttt{points(fwd, spot, pip)}~;
\texttt{implied\_rate\_quote}~; \texttt{basis\_on\_quote}~; \texttt{fx\_swap\_legs}~;
\texttt{hedge\_cost}~; \texttt{hedged\_yield}.
\bfield{Règles} Taux simples du marché monétaire avec la convention de décompte de chaque devise~;
base sur la jambe hors dollar~; \omterm{def:m2:fx-swaps-forwards-and-the-cross-currency-basis:forward}{points de terme} exprimés dans les pips de la paire.
\bfield{Tests d'acceptation} \texttt{code/firm/fxfwd/tests/}~: CIP du terme et son inverse~; la base
retrouvée à partir d'un terme construit avec elle~; les jambes d'un \omterm{def:m2:fx-swaps-forwards-and-the-cross-currency-basis:swap}{swap de change} qui se compensent
exactement à l'écart de taux d'intérêt~; le rendement couvert plus faible avec une base négative.
\bfield{Pour aller plus loin} Courbes de termes à partir des courbes d'actualisation complètes dans
chaque devise (\cref{bld:m2:interest-rate-swaps:curve})~; dates irrégulières~; un \omterm{def:m2:fx-swaps-forwards-and-the-cross-currency-basis:basis}{swap de base de devises}
évalué sur deux courbes, la base jouant le rôle d'un spread.
\end{build}

\begin{omsources}
\item C.~Borio, R.~McCauley, P.~McGuire et V.~Sushko, ``Covered interest parity lost: understanding the cross-currency basis'', \emph{BIS Quarterly Review},
septembre 2016.
\item Banque des règlements internationaux, Triennial Central Bank Survey 2025.
\item Federal Reserve, communiqués de presse des 12 décembre 2007, 13 octobre 2008, 15
et 19 mars 2020~; publications H.4.1 et H.10.
\item Séries FRED DGS10, SOFR, IRSTCI01JPM156N, IRLTLT01JPM156N et SWPT.
\end{omsources}

\section{Exercices}

\begin{exercise}[$\star$]\label{exo:m2:fx-swaps-forwards-and-the-cross-currency-basis:1}
EURUSD vaut 1.1464, le taux en euros 2.00~\% et le taux en dollars 3.68~\%, tous deux
en actual/360. Calculer le terme à 90 jours et ses points.
\end{exercise}

\begin{exercise}[$\star$]\label{exo:m2:fx-swaps-forwards-and-the-cross-currency-basis:2}
Pourquoi la devise au taux d'intérêt le plus élevé se négocie-t-elle à un terme en déport~?
\end{exercise}

\begin{exercise}[$\star$]\label{exo:m2:fx-swaps-forwards-and-the-cross-currency-basis:3}
Dans le swap de la \cref{fig:m2:fx-swaps-forwards-and-the-cross-currency-basis:legs}, qui a emprunté
quelle devise, et quel est le coût en yens des dollars~?
\end{exercise}

\begin{exercise}[$\star\star$]\label{exo:m2:fx-swaps-forwards-and-the-cross-currency-basis:4}
Le terme USDJPY à trois mois du marché vaut 155.7059 avec les taux de
l'\cref{ex:m2:fx-swaps-forwards-and-the-cross-currency-basis:fwd}. Quelle base
ce terme implique-t-il, et que signifie son signe~?
\end{exercise}

\begin{exercise}[$\star\star$]\label{exo:m2:fx-swaps-forwards-and-the-cross-currency-basis:5}
Un trader long de USD~100 millions contre yens reporte sa position avec un \omterm{def:m2:fx-swaps-forwards-and-the-cross-currency-basis:swap}{tom-next}.
En utilisant les points \omterm{def:m2:fx-swaps-forwards-and-the-cross-currency-basis:swap}{tom-next} de l'\cref{ex:m2:fx-swaps-forwards-and-the-cross-currency-basis:fwd},
combien un report rapporte-t-il ou coûte-t-il, en yens~?
\end{exercise}

\begin{exercise}[$\star\star$]\label{exo:m2:fx-swaps-forwards-and-the-cross-currency-basis:6}
Pourquoi les fins de trimestre font-elles bouger la base, et pourquoi les lignes de swap de mars 2020
l'ont-elles réduite~?
\end{exercise}

\begin{exercise}[$\star\star\star$]\label{exo:m2:fx-swaps-forwards-and-the-cross-currency-basis:7}
\emph{Programmation.} Avec \texttt{hedged\_series}, trouver le premier mois des données
où le Treasury couvert à base nulle rapporte moins que le JGB, et compter de tels
mois.
\end{exercise}

\begin{exercise}[$\star\star\star$]\label{exo:m2:fx-swaps-forwards-and-the-cross-currency-basis:8}
\emph{Repérer l'erreur.} «~La base est négative, donc une banque peut emprunter des dollars sur
le marché monétaire, les prêter par des \omterm{def:m2:fx-swaps-forwards-and-the-cross-currency-basis:swap}{swaps de change} et verrouiller un profit sans risque~;
cela fera disparaître la base en quelques jours.~» Corriger cette affirmation.
\end{exercise}

\section{Problème~: le rendement couvert}

\begin{problem}[{Problème du week-end --- ce que rapporte un Treasury à un investisseur en yens}]\label{pb:m2:fx-swaps-forwards-and-the-cross-currency-basis:1}
Un assureur-vie japonais envisage d'acheter des Treasuries à dix ans à 4.75~\% le 31
août 2026 et de les couvrir en yens par des \omterm{def:m2:fx-swaps-forwards-and-the-cross-currency-basis:swap}{swaps de change} à trois mois renouvelés. Les
taux en dollars sont de 3.68~\%, les taux en yens de 0.977~\%, le cours USDJPY à 156.87~; le JGB à dix ans rapporte 2.94~\%.

\medskip\noindent\textbf{Partie I --- La couverture.}
\begin{enumerate}
\item Donner le terme à trois mois sans base, et ses points.
\item Donner le coût annualisé de la couverture sans base.
\item Donner le rendement couvert sans base.
\item Le donner avec une base de $-25$ puis de $-50$ points de base.
\item Comparer chacun de ces rendements au JGB.
\end{enumerate}

\medskip\noindent\textbf{Partie II --- La mécanique.}
\begin{enumerate}[resume]
\item Que fait l'assureur à chaque renouvellement~?
\item Que se passe-t-il pour les flux de trésorerie de la couverture si le dollar baisse de 10~\% en un trimestre~?
\item Pourquoi l'assureur a-t-il besoin de liquidité en yens, et quand~?
\item En quoi une couverture à terme d'un an diffère-t-elle de swaps à trois mois renouvelés~?
\item Qu'est-ce qui rendrait la couverture moins coûteuse~?
\end{enumerate}

\medskip\noindent\textbf{Partie III --- Les alternatives.}
\begin{enumerate}[resume]
\item Que gagne l'assureur sans couverture, et que risque-t-il~?
\item Quel seuil de variation du dollar faut-il sur un an pour que la position non couverte batte le JGB~?
\item Pourquoi pourrait-il ne couvrir que la moitié~?
\item Pourquoi de tels flux élargissent-ils la base~?
\item Qui prend la contrepartie de ses \omterm{def:m2:fx-swaps-forwards-and-the-cross-currency-basis:swap}{swaps de change}~?
\end{enumerate}

\medskip\noindent\textbf{Partie IV --- Jugement.}
\begin{enumerate}[resume]
\item Quand les Treasuries couverts ont-ils battu pour la dernière fois les JGB dans les données, et pourquoi à ce moment-là~?
\item Comment une baisse de 100 points de base du taux de la Federal Reserve modifierait-elle la réponse~?
\item Pourquoi la base que paie l'assureur n'est-elle pas un repas gratuit pour la banque qui prête les dollars~?
\item Énoncer le \emph{résultat nommé}~: le rendement couvert du Treasury pour l'assureur.
\item En une phrase~: que gagne un investisseur couvert en obligations étrangères~?
\end{enumerate}
\end{problem}

\section{Questions d'entretien}

\begin{interviewq}[$\star$ \iqroles{trader, chercheur}]\label{iq:m2:fx-swaps-forwards-and-the-cross-currency-basis:1}
Démontrez le taux de change à terme à partir du taux au comptant et des taux d'intérêt.
\end{interviewq}

\begin{interviewq}[$\star$ \iqroles{banque, trader}]\label{iq:m2:fx-swaps-forwards-and-the-cross-currency-basis:2}
Qu'est-ce qu'un \omterm{def:m2:fx-swaps-forwards-and-the-cross-currency-basis:swap}{swap de change}, et pourquoi est-ce l'instrument de change le plus échangé~?
\end{interviewq}

\begin{interviewq}[$\star\star$ \iqroles{chercheur, trader}]\label{iq:m2:fx-swaps-forwards-and-the-cross-currency-basis:3}
Pourquoi la \omterm{def:m2:fx-swaps-forwards-and-the-cross-currency-basis:cip}{parité des taux d'intérêt couverte} ne tient-elle plus depuis 2008~?
\end{interviewq}

\begin{interviewq}[$\star\star$ \iqroles{trader}]\label{iq:m2:fx-swaps-forwards-and-the-cross-currency-basis:4}
Comment couvririez-vous un portefeuille d'obligations en euros pour un investisseur en dollars, et que
rapporte ou coûte la couverture~?
\end{interviewq}

\begin{interviewq}[$\star\star$ \iqroles{banque}]\label{iq:m2:fx-swaps-forwards-and-the-cross-currency-basis:5}
Que se passe-t-il sur le marché des \omterm{def:m2:fx-swaps-forwards-and-the-cross-currency-basis:swap}{swaps de change} à la fin d'un trimestre, et que ferait une banque
à ce sujet~?
\end{interviewq}

\begin{interviewq}[$\star\star\star$ \iqroles{développeur, chercheur}]\label{iq:m2:fx-swaps-forwards-and-the-cross-currency-basis:6}
Concevez un service qui publie les \omterm{def:m2:fx-swaps-forwards-and-the-cross-currency-basis:forward}{points de terme} pour cinquante paires et toutes les
échéances standard, ainsi que la base implicite, à partir des courbes de taux et des cotations de marché.
\end{interviewq}
